\subsection{The two moment tests on the decorated lattices}
\label{subsec:word-moments}

To apply the assigned-prime bound to the selected columns, express the kernel at the decorated lattices through its two moment tests.

At a standard pair at \(T\), let \(D_v,D'_v\) be the ordinary
local lattices with residue rows \(\ell_v,\ell'_v\). The word
character is unramified on the compact determinants there.
The scalar of the kernel in the standard frames is a unit:
the determinant ideals have exponent zero at the places over
\(T\), the normalized away characters are units there, and
the finite level characters have finite image. The two moment tests imply \eqref{eq:tag-w1} with its distinguished-place loss. Write \(T=p\) in this subsection and set
\(F=fN_i\).

For a rational matrix \(x\) in the Hom space put \(a_{\tau,ij}(x)=(u_\tau\lambda(\tau(x^*)))_{ij}\), with \(u_\tau\) the prescribed unit. Its point evaluation at the multi-index \(m\) is
\[
  \prod_{\tau,i,j}a_{\tau,ij}(x)^{m_{\tau,ij}}.
\]
 Let \(J_m\) be the Taylor jet with the
derivation generators of the stage and its finite reader.
Retain the indices whose determinant-normalized
jet at each contributing coset belongs to the coefficient field.
These are the rational indices used by the stage contraction.
Sections~\ref{subsec:word-measure} and~\ref{subsec:word-projection}
will construct the two approximations needed in the next statement.
We first show why their different supports are sufficient.

\begin{proposition}[A moment criterion for the word prime]\label{prop:word-moment-criterion}
Use the selected rational jet indices on the standard lattices at
\(p=T\). At every prescribed precision and for every finite list
of these indices, suppose that the jet vector can be approximated
by one integral linear combination of point evaluations times
\(p^{-F}\) at the distinguished place, and by one integral linear
combination without that factor at the conjugate place. The
respective supports of the evaluations must satisfy
\[
\begin{aligned}
 &x^*(\ker\ell'_v)\subset D_v\quad\text{for every }v,\\
 &xD_v\subset D'_v,\qquad \ell'_v x=\ell_v
                      \quad\text{on }D_v\text{ for every }v .
\end{aligned}                                                     \tag{W6}\label{eq:tag-w6}
\]
With the selected target weights in the lowest alcove, these
approximations imply the two pointwise bounds \eqref{eq:tag-w1}. The same
implication holds for the transported decorations at \(p\).
\end{proposition}

\begin{proof}
\proofstep{Identify the conjugate lattice map.}
At a conjugate coefficient place the adjoint uses the opposite
local factors. Write \(O_v\) for the ring of integers in the local row
coordinates. With integral lifts of the residue rows, their kernel duals are
\[
 (\ker\ell_v)^\vee=D_v^\vee+O_v\,p^{-1}\ell_v,\qquad
 (\ker\ell'_v)^\vee=(D'_v)^\vee+O_v\,p^{-1}\ell'_v
\]
in the local residue-row coordinates. The formula follows by
taking a basis in which the row modulo \(p\) is the first
coordinate. By \eqref{eq:tag-w6}, \(x^*(D'_v)^\vee\subset D_v^\vee\) and
\[
 x^*(p^{-1}\ell'_v)-p^{-1}\ell_v
    =p^{-1}(\ell'_v x-\ell_v)\in D_v^\vee .
\]
Together these inclusions give the adjoint map on the two
enlarged dual lattices.

\proofstep{Identify the integral polynomial lattice at the target.}
For a homogeneous source polynomial, contraction against the
monomials of a single point evaluation is polynomial substitution
by its adjoint matrix, with the prescribed determinant shifts.
This follows on every monomial by the multinomial theorem for
the substitution, or by applying the commuting derivations to
the exponential with this linear argument. In adapted bases the
matrix is integral by \eqref{eq:tag-w6}, so Cauchy--Binet gives integrality
of the polynomial part. The source Weyl lattice is contained in
the dual Weyl lattice of polynomials. At the target these lattices
coincide over the local ring.

To see this, if the target has rank
\(r\) and dominant weight \(\nu\), its chosen size satisfies
\[
     \nu_a-\nu_b+b-a\le \nu_1-\nu_r+r < p
                    \qquad (1\le a<b\le r).
\]
After the determinant twist the weight is polynomial. In
characteristic \(p\) the Jantzen sum for its Weyl module is
empty: the multiples \(mp\) in the formula satisfy
\(0<mp<\langle\nu+\rho,\alpha^\vee\rangle\).
Thus the Weyl module and its contravariant dual are simple.
The standard map from the integral Weyl lattice to the dual
Weyl lattice sends a highest generator to the primitive highest
minor. Its reduction therefore has full image over an algebraic
closure, and hence over the residue field. Nakayama's lemma gives
the integral equality, also after each change of adapted basis.
\cite{jantzen-rep}\label{cor:word-moments-1}

\proofstep{Pass to the jet limit and account for determinant shifts.}
At the prescribed degree write \(C\) for the finite linear
map from a vector of moments to the contracted coefficients.
For each precision the hypotheses supply a vector \(u_n\) of
approximating moments on the required support, including its
branch's factor in \eqref{eq:tag-w6}. Then \(u_n\to (J_m)_m\), so

\[
  C u_n\to C(J_m)_m.
\]
 All \(C u_n\) belong to the same
rescaled integral lattice by the pointwise contraction and
determinant bounds. This lattice is closed, so the limit does
too. Coefficients in the chosen number field recover their
integrality by restricting the valuation along its embedding. At the
distinguished place choose the identity at the source in the
standard basis and a basis of the kernel of the target row.
Its determinant can be taken equal to \(p\), since the row is
nonzero modulo \(p\). The contribution to the target determinant
factor is 
\[
  p^{\sum_\tau\nu_{{\rm last},\tau}}
=p^{-\sum_\tau s_{{\rm div},i}(\tau)}.
\]

At the conjugate place the two enlarged dual lattices have
basis determinants of valuation \(-1\) in each position. For
source and target shifts \(\sigma,\sigma'\), the determinant
factor on substitution is \((\det B)^{\sigma'}
(\det B')^{-\sigma}\) per position. The old coordinates of
the stage weights give \(\sigma-\sigma'=(t+1)/2=s\).
Thus at each position the product has valuation
\(\sigma'(-1)-\sigma(-1)=\sigma-\sigma'=s\).
Over the \(f\) positions the contributions sum to
\[
       v_p\!\left(\prod_\tau(\det B_\tau)^{\sigma'_\tau}
                                 (\det B'_\tau)^{-\sigma_\tau}\right)
                       = fs.
\] Together with the
two powers multiplying the point evaluations these are exactly
\eqref{eq:tag-w1}.

\proofstep{Transport the support and the jet to other decorations.}
For nonstandard decorations choose rational unitary changes
taking the lattices to standard ones above \(p\). Such changes
exist by weak approximation: in a dense Cayley chart on the
unitary group the inverse image of the prescribed open sets of
lattice-carrying matrices is open in a rational vector space.
At a split place a general linear change and its inverse adjoint
give the desired target, and one can choose it in the chart by
openness and density. Transport the finite away reader and the
frames by the same rational maps. They take both the value
lattices and the source and target rows to their transported
versions. The change-of-variables identity for the rational
kernel carries \eqref{eq:tag-w1} back.
\end{proof}

Sections~\ref{subsec:word-measure} and~\ref{subsec:word-projection}
prove the simultaneous approximations on these two supports
for the selected normalized vector. Thus \eqref{eq:tag-w1} holds on
the stage readings for Proposition~\ref{prop:word-prime-transfer}, and \eqref{eq:tag-w2}--\eqref{eq:tag-w5} give the sharp
contribution at the prime assigned to the sequence on the trace-zero column.
